\section{The abelian cases of the index route}
\label{sec:abelian}

We use $C_{\mathrm G}(X)$ for the historical condition on cote~16:
the normalized lowering operator $\Lambda_X$ is algebraic.  This is
$B(X)$ in \cite[\S2]{Kleiman1968}.  It is stronger than the algebraicity
of the K\"unneth projectors, now usually called $C(X)$.
Throughout, an algebraic class means a rational cycle class modulo the
specified Weil homological equivalence.  The statements do not replace
homological equivalence by rational or algebraic equivalence.

\begin{proposition}[SN-10: explicit critical inverse maps in the abelian cases]
Let $A$ be an abelian variety of dimension $g$ over an algebraically closed
field, let $h$ be an ample rational divisor class, and use Betti cohomology
for a complex variety (or after spreading and embedding in $\mathbb C$),
or $\ell$-adic cohomology with $\ell$ invertible.
Then $C_{\mathrm G}(A)$ holds.  If $j:Y\hookrightarrow A$ is a smooth
hyperplane section for the chosen polarization, the following maps are
algebraic:
\[
 s=\Lambda_A j_*:H^k(Y)\longrightarrow H^k(A),\qquad
 t=j^*\Lambda_A:H^{k+2}(A)\longrightarrow H^k(Y).
\]
They satisfy $sj^*=1$ for $k\leq g-1$ and $j_*t=1$ for $k\geq g-1$.
The same statements apply to every power and product of abelian varieties.
Consequently the correctly indexed weak algebraicity conditions, as well
as the weighted critical Gysin conditions used in the corrected square
criterion, hold in these cases.  Moreover $C_{\mathrm G}(Y)$ holds.
\end{proposition}
\begin{proof}
The algebraicity assertion is the classical abelian Lefschetz theorem,
\cite[Theorem~2A11]{Kleiman1968}; in the stronger divisor-generated form it
is \cite[Theorem~5.9, with the author's correction to its proof]{Milne1999a}.
We cite that result rather than reprove it.  Here $\Lambda$ means unit
lowering on a Lefschetz string:
\[
 \Lambda(L^r v)=L^{r-1}v\ (r>0),\qquad \Lambda v=0
 \quad(v\text{ primitive}).
\]
No string terminates in degree less than $g$ and no string starts in degree
greater than $g$.  Therefore $\Lambda L=1$ on $H^k(A)$ for $k<g$, and
$L\Lambda=1$ on $H^{k+2}(A)$ for $k+2>g$.  The projection formula
$j_*j^*=L$ proves the two asserted identities, including the endpoints.
All factors are algebraic correspondences, so the identities give the
claimed descent and lifting of algebraic classes.  Explicitly, if $T$ has
even dimension $d$ and $j:W\hookrightarrow T$, a right inverse to
$L_Tj_*:H^{d-2}(W)\to H^{d+2}(T)$ is $j^*\Lambda_T^2$:
$L_Tj_*j^*\Lambda_T^2=L_T^2\Lambda_T^2=1$ in degree $d+2$.
Every Lefschetz string occurring there has at least two preceding steps,
which verifies the last equality.  Products remain abelian.  Finally,
$\Lambda_Y=j^*\Lambda_A^2j_*$ by cote~16 scan~040, equivalently
\cite[Proposition~2.12]{Kleiman1968}; this proves the last assertion.
\end{proof}

\begin{theorem}[SN-11: completion of the abelian index route]
Under the supplied CM Hodge theorem, every abelian variety $A$ over an
algebraically closed field satisfies $C_{\mathrm G}(A)$ and $I(A\times A)$.
In particular, for every rational ample $h$ on $A$ and every $2r\leq g$,
\[
 (x,y)\longmapsto (-1)^r\deg(xy h^{g-2r})
 \quad\text{on }\ker(h^{g-2r+1}:A^r_{\rm hom}(A)\to A^{g-r+1}_{\rm hom}(A))
\]
is positive definite over $\mathbb R$; homological and numerical
equivalence coincide.  The rational algebra of algebraic correspondences
modulo homological equivalence is semisimple, and its subalgebras stable
under the Lefschetz adjoint are semisimple.  Algebraic injections and
surjections between the cohomological summands of abelian varieties split
by algebraic correspondences.

For powers of a CM elliptic curve, including the hexagonal curve of
HODGE\_SIGN SG-04, the same conclusions use SG-04 and Milne's positivity
for Lefschetz classes and do not require the supplied general CM Hodge
theorem.
\end{theorem}
\begin{proof}
The first input is precisely the supplied 032, Theorem~1.1 (p.~1 and its
concluding proof on p.~47), equivalently van Hoften--Sempliner,
Hexagon~2610.00259v2, Theorem~I, p.~1.  Milne's published
Theorem~3.3(a),(b) gives the asserted equality of equivalence relations and
the index inequality in every characteristic; the attached 032,
Corollary~8.4, p.~50, states exactly this consequence, including rational
ample polarizations.  Apply it to $A\times A$ with $h\boxtimes1+1\boxtimes h$.
SN-10 supplies the other hypothesis of the classical positive-adjoint
theorem \cite[Theorem~3.11]{Kleiman1968}.  Its Corollaries~3.12 and~3.14
then give semisimplicity and splitting.  These are cited consequences,
not new proofs of those classical theorems.  All products of the abelian
varieties in a proposed correspondence are again abelian, so the product
positivity hypotheses in those corollaries are satisfied.

For elliptic powers, SG-04 identifies all algebraic class spaces with the
divisor-generated spaces, on the relevant complex varieties and their
good reductions.  \cite[Remarks~3.7--3.8]{Milne2002} supplies positivity on
these spaces.  The same application of SN-10 and Kleiman's corollaries
therefore applies.  This last argument cites SG-04 and does not reprove it.
\end{proof}

\paragraph{Exact disposition of Q2c.}
The intended positivity question Q2c is closed by the preceding cited
chain: 032 Theorem~1.1 $\Rightarrow$ Milne~2002 Theorem~3.3 $\Rightarrow$
032 Corollary~8.4.  The human deposit's \S1.1.2 uses
Theorem~4.6(a),(b), the numbering of Milne's original version
(the author's \texttt{2002aS.pdf}).  The later version corresponding to
the published paper has the needed conclusion as Theorem~3.3(a),(b)
(\texttt{2002aP.pdf}, pp.~8--9).  Both versioned locators have been
checked; this is a version crosswalk, not a mathematical discrepancy.

\subsection{A register correction: algebraic does not mean divisor-generated}

\begin{proposition}[SN-12: a concrete counterexample to the literal U-02 wording]
With the supplied 032 theorem, there is a simple complex CM abelian
fourfold with algebraic codimension-two classes that are not generated by
divisors.  Thus the predecessor's literal U-02 sentence
``Every algebraic class on an arbitrary CM abelian variety is a Lefschetz
class'' is false, even though the intended Q2c positivity question is
closed by SN-11.
\end{proposition}
\begin{proof}
Let $F=\mathbb Q(\alpha)$, where
\[
 f(T)=T^4-5T^2+T+1,\qquad K=\mathbb Q(i),\qquad E=FK.
\]
The signs of $f$ give four distinct real roots, respectively in
$(-3,-2),(-1,0),(0,1),(2,3)$.
Modulo $3$, $f=T^4+T^2+T+1$ has no linear root.  Its remainders modulo the
three monic irreducible quadratics over $\mathbb F_3$,
$T^2+1,T^2+T+2,T^2+2T+2$, are respectively
$T+1,1,2T+1$.  It is therefore irreducible modulo $3$.
Modulo $2$, its factorization is
$(T+1)(T^3+T^2+1)$, with the cubic irreducible and the factors distinct.
The Galois group of its splitting field $N$ is thus a transitive subgroup
of $S_4$ containing a $4$-cycle and a $3$-cycle, hence $S_4$.
The field $N$ is totally real, so $N\cap K=\mathbb Q$.

The only proper CM subfield of $E$ is $K$.  Here is the subgroup check,
to make this assertion explicit.  In
$\operatorname{Gal}(NK/\mathbb Q)=S_4\times C_2$ the subgroup fixing $E$
is $S_3\times\{1\}$.  A subgroup containing it projects to either $S_3$
or $S_4$.  In the first case it is $S_3\times\{1\}$ or $S_3\times C_2$.
In the second case it is the full group or the graph of a homomorphism
$S_4\to C_2$ trivial on $S_3$; that homomorphism is trivial, so the subgroup
is $S_4\times\{1\}$.  The intermediate fields are therefore just
$E,F,K,\mathbb Q$.

Choose a CM type of $E$ with two embeddings above each embedding of $K$.
It cannot be induced from $K$, so it is primitive.  The CM classification
\cite[Propositions~3.12--3.13]{MilneCM} produces a simple abelian fourfold
$A$ with $\operatorname{End}^0(A)=E$ of this type.  The action of $K$ on
$H^{1,0}(A)$ has signature $(2,2)$.  Hence the rational two-dimensional
Weil subspace
\[
 W=\bigwedge_K^4 H^1(A,\mathbb Q)\ \subset\ H^4(A,\mathbb Q)
\]
has Hodge type $(2,2)$.  The inclusion here is the canonical determinant
subspace: after splitting $K$, its two lines are
$\bigwedge^4 V_+$ and $\bigwedge^4 V_-$ for the two $K$-eigenspaces in
$H^1(A,\mathbb C)$.

Rosati on the CM field $E$ is complex conjugation.  The
N\'eron--Severi/Rosati identification
\cite[Proposition~13.1]{MilneAV} gives
$\dim_{\mathbb Q}\operatorname{NS}(A)_{\mathbb Q}=\dim_{\mathbb Q}F=4$.
In the decomposition of $H^1$ into the eight $E$-embedding lines,
the four divisor directions pair an embedding with its conjugate.
For clarity, this last description does not follow merely by a dimension
count: a polarization pairs only conjugate lines, because
$\psi(ax,y)=\psi(x,\bar a y)$, and multiplying it by $F$ gives the four
independent such pairs.  Every product of two divisor classes consequently
has $K$-bidegree $(2,2)$.  The two lines of $W$ have bidegrees $(4,0)$ and
$(0,4)$.  Their intersection with the divisor-generated subspace is zero.
Every nonzero rational element of $W$ is algebraic by 032 Theorem~1.1.
This gives the asserted counterexample.
\end{proof}

\section{The additional inputs for the positive-adjoint argument}
\label{sec:positive}

Let $R_X$ denote the image of rational algebraic correspondences in
$\operatorname{End}(H^*(X))$.  It is the image algebra, so classes acting
as zero have already been quotiented out.  Under $C_{\mathrm G}(X)$,
the Lefschetz star and degree projectors are algebraic.  The adjoint
$u\mapsto u^\dagger$ with respect to the Lefschetz pairings preserves
$R_X$, and
\[
  \tau(u)=\sum_i\operatorname{Tr}(u_{ii}\mid H^i(X))\in\mathbb Q.
\]
The sum here is the ordinary trace, not the alternating trace.  The degree
projectors convert the cohomological Lefschetz trace formula to this
rational trace.  The extra condition needed by scan~049 is exactly
\begin{equation}\label{eq:positive-trace}
             \tau(u^\dagger u)>0\qquad(0\ne u\in R_X).
\end{equation}
More invariantly, the associated symmetric form on $R_X\otimes\mathbb R$
must be positive definite.  Kleiman~1968, Theorem~3.11, derives precisely
this from $I(X\times X)$ with the product polarization.  Positivity on
only divisor classes or on only the scalar Frobenius traces is not this
condition.

In an $\ell$-adic realization, this positivity also supplies the
finite-dimensionality over $\mathbb Q$ needed below. Indeed, if
$u_1,\ldots,u_s$ are rationally independent in $R_X$, their rational
Gram matrix $\bigl(\tau(u_i^\dagger u_j)\bigr)$ is positive definite
and hence nonsingular. Any $\mathbb Q_\ell$-linear relation among
their cohomological operators, paired with each $u_i^\dagger$, is
therefore zero only with all coefficients zero. Thus
$s\le(\dim_{\mathbb Q_\ell}H^*(X))^2$. This avoids inferring
rational finite-dimensionality merely from an inclusion in a
finite-dimensional $\mathbb Q_\ell$-space.

\begin{proposition}[SN-13: constructive splitting from the positive trace form]
Let $R$ be a finite-dimensional rational algebra acting faithfully on
a vector space over a characteristic-zero field, equipped with an
involution $\dagger$ and a cyclic rational trace $\tau$ satisfying
$\tau(a^\dagger)=\tau(a)$ and such that
$\tau(a^\dagger b)$ is positive definite over $\mathbb R$.
Let $e,f$ be self-adjoint idempotents
and $u=fue$.  If $u:eV\to fV$ is injective, then
\[
 v=(u^\dagger u)^{-1}u^\dagger\in eRf,
 \qquad vu=e,
\]
where the inverse is taken in $eRe$.  If $u$ is surjective, the analogous
right inverse is $u^\dagger(uu^\dagger)^{-1}$.
\end{proposition}
\begin{proof}
On $R_{\mathbb R}$ use the inner product
$\langle a,b\rangle=\tau(a^\dagger b)$.
It is symmetric because $\tau(a^\dagger)=\tau(a)$, and is positive
definite by hypothesis; both hold for the geometric trace form above.
Left multiplication by $u^\dagger$ is the adjoint of left multiplication
by $u$.  If $a\in eR$ and $ua=0$, faithfulness of the action on $V$ and
injectivity on $eV$ imply $a=0$.  Thus $u^\dagger u$ acts invertibly on
$eR$, since
$\langle a,u^\dagger ua\rangle=\langle ua,ua\rangle>0$ for $a\ne0$.
Its minimal polynomial on this finite-dimensional rational vector space
has nonzero constant coefficient.  Bezout with $T$ gives an inverse which
is a rational polynomial in $u^\dagger u$ with identity $e$, hence belongs
to $eRe$.  Multiplication verifies $vu=e$.  For a surjection,
$a\mapsto au$ is injective on $Rf$: if $au=0$, then $a$ vanishes on
$fV=\operatorname{im}(u)$ and $a=af$, so $a=0$.
Cyclicity shows that right multiplication by $u^\dagger$ is the adjoint
of right multiplication by $u$.  Thus right multiplication by
$uu^\dagger$ is invertible on $Rf$.  The same polynomial argument gives
its inverse in $fRf$, and multiplication verifies the stated right inverse.
\end{proof}

\begin{theorem}[SN-14: trace and $\lambda$ identities do not supply scan-049 positivity]
Even if all coefficients from a specified correspondence algebra are
allowed in traces, pure Frobenius eigenvalues, effective Grothendieck-group
identities, and their tensor, exterior-power and dual consequences do not
formally imply semisimplicity of that algebra or
\eqref{eq:positive-trace}.  They also do not imply semisimplicity of a
pure sheaf's Frobenius at a point.  These are counterexamples to a formal
implication, not counterexamples to the standard conjectures for varieties.
\end{theorem}
\begin{proof}
Work in the rigid tensor category of rational representations of
$\mathbb G_a\times\mathbb G_m$.  Take $V=\mathbb Q^3$ and
\[
 N=\begin{pmatrix}0&1&0\\0&0&1\\0&0&0\end{pmatrix},\qquad
 J=\begin{pmatrix}0&0&1\\0&-1&0\\1&0&0\end{pmatrix},
\]
with action $(t,s)\mapsto s\exp(tN)$.  Its endomorphism algebra is
$R=\mathbb Q[N]\simeq\mathbb Q[\epsilon]/(\epsilon^3)$:
an endomorphism commutes with the action exactly when it commutes with
the single Jordan block $N$, whose centralizer consists of its
polynomials.  Put $F=c\exp(N)$ with $c=q^r>0$ for an integer $r\geq0$.
All Frobenius eigenvalues have pure weight $2r$, while both $F$ and
the algebra $R$ are nonsemisimple.  The $J$-adjoint satisfies
$N^\dagger=-N$, so $R$ is closed under it and $F^\dagger F=c^2I$.
Nevertheless $\tau(N^\dagger N)=-\operatorname{Tr}(N^2)=0$.
For every $a=a_0I+a_1N+a_2N^2$ and every integer $s\geq0$,
\[
 \operatorname{Tr}(aF^s)=3a_0c^s.
\]
Thus all these coefficient traces factor through $R/\operatorname{rad}R$.

More generally the invariant filtration
$0\subset\ker N\subset\ker N^2\subset V$
has three one-dimensional trivial unipotent quotients.  It is compatible with
$F$, and the standard filtrations of tensor products, exterior powers,
and duals have associated graded built from the same quotients.  Traces
are additive in these filtrations.  Consequently all the indicated
character identities, including insertions of the induced coefficient
operators, are identical to those on the associated graded.
In the Grothendieck $\lambda$-ring of this tensor category this is simply
$[V]=3[\mathbb Q_c]$ and
$\lambda_t([V])=(1+[\mathbb Q_c]t)^3$.
The notation $\mathbb Q_c$ denotes the one-dimensional torus character
evaluated at $F$, where it has eigenvalue $c$.
Even effectivity does not help: $V$ itself is an actual object.

For the sheaf assertion, on $\operatorname{Spec}\mathbb F_q$ take the
continuous $\ell$-adic representation of $\widehat{\mathbb Z}$ defined by
$z\mapsto\exp(z_\ell N)$, where $z_\ell$ is its projection to
$\mathbb Z_\ell$.  This is a continuous homomorphism into
$\operatorname{GL}_3(\mathbb Q_\ell)$, since the exponential is the finite
polynomial $I+z_\ell N+z_\ell^2N^2/2$.
Its Frobenius is $\exp(N)$ (or its inverse under the other
Frobenius convention).  At every extension point the eigenvalues are all
$1$, so this is a punctually pure sheaf of weight zero.  It is not
semisimple as a Frobenius representation and has
$\det(1-TF)=(1-T)^3$, identical to the split sheaf.
\end{proof}

\paragraph{Application to the sign-only theorem.}
HODGE\_SIGN SG-08--SG-11, together with GCT\_NSRH F03, is used as the
settled sign-only theorem stipulated in the prompt.  Its conclusion and
effective remainder are not reproved.  Applying its trace identities to
a product changes neither the preceding filtration argument nor the
distinction between a correspondence and a lisse coefficient sheaf.
A correspondence acting on $H^*(X)$ does not automatically preserve a
chosen generic hyperplane and its Gysin kernel.  Even when such
functorial compatibility is furnished, the trace-only information in
SN-14 remains insensitive to the radical.  To complete the
scan-049/Kleiman argument by this method, one must additionally establish
the algebraic adjoint and the nondegenerate positive form
\eqref{eq:positive-trace}; asserting either from the parity of an
effective remainder would assume the missing conclusion.  SN-14 does
not exclude a different semisimplicity argument.

\paragraph{Motivic and arithmetic formulations.}
For comparison, numerical motives with rational coefficients are already
a semisimple abelian category by Jannsen,
\emph{Motives, numerical equivalence, and semi-simplicity},
Theorem~1 \cite{Jannsen1992}.  That cited theorem concerns numerical
equivalence.  It does not identify numerical equivalence with the
homological equivalence used by the given realization, nor does it
provide the positive Lefschetz-adjoint form on that realization.
Those are the extra assertions supplied in the abelian case here.

SN-11 gives the semisimple polarized category of abelian motives modulo
$\ell$-adic homological equivalence, identified with its numerical
category on all products.  Over a finite field the Lefschetz adjoint makes
normalized Frobenius unitary, so it is semisimple; the classical exact
locator is \cite[Lemma~4.6 and Theorem~4.7]{Kleiman1968}.
The absolute values already come from the supplied sign-only theorem.
Splittings of motives split their local characteristic polynomials and
therefore multiply their local $L$-factors.  SN-14 explains why the
converse from these factors or their $K_0$ classes fails.

For a finite collection of the varieties, polarizations, cycles and
operators in SN-10--SN-13, choose a finitely generated ring over
$\mathbb Z$ carrying the data and shrink to a smooth projective model;
finite-presentation descent is EGA~IV$_3$, 8.8.2, and smooth proper
cohomological specialization is SGA~4, XVI, 2.2.
The realization consists of the lisse sheaves $R^i f_*\mathbb Q_\ell$ and
the corresponding morphisms of sheaves.  The trace/correspondence
operations are SGA~5, III, 3 and 5.1--5.3; the $\lambda$-ring identity in
SN-14 uses the exact-sequence rule SGA~6, V, 2.1--2.2, especially
Lemma~2.2.1.  A $K_0$ class
records that rule but does not carry a polarization or an integral
splitting.  This identifies the precise place of these findings in
cote~16 scans~040, 049--050 and in the predecessor's U-01/U-02, without
promoting them to claims about Chow groups.
